\subsection{Wysp Core}

\subsubsection{Modules}

\begin{xltabular}{\textwidth}{@{}L{2.7cm}Y@{}}
\toprule
\textbf{Module} & \textbf{Responsabilité} \\
\midrule
\file{contract/} & Modèles Pydantic du contrat et du protocole du panneau, export des schémas, prédicats, empreinte d'état. Seul module importé par les SDK et le kit. \\
\file{gateway/} & Serveur JSON-RPC \code{/ws/adapter}, \code{AdapterProxy} typé avec budgets de temps, registre des adaptateurs, surveillance des battements, contrôle d'origine et de jeton. \\
\file{session/} & Session de tutorat (adaptateur actif, profil, tableau, narration, propositions), bus d'événements asyncio, mémoire de conversation avec résumé glissant, constructeur de contexte. \\
\file{orchestrator/} & Boucle de tour (\code{TurnRunner}), outils offerts au LLM, prompts génériques, annulation. \\
\file{llm/} & Interface \code{LLMProvider}, \code{ScriptedLLM} (scénarios rejouables), fournisseur compatible OpenAI (OpenAI, Groq, Ollama local…), Bedrock en option. \\
\file{board/} & Modèle visuel, opérations et correctifs versionnés, projection des entités, validateur de provenance, regroupement au-delà de 60 nœuds. \\
\file{narration/} & Segments, repères de synchronisation, curseur de lecture, interruption et reprise. \\
\file{trust/} & Politique de risque, porte de permission, autorisations de leçon, pile d'annulation, journal d'audit. \\
\file{verification/} & Observation fraîche, évaluation des prédicats, preuves. \\
\file{pedagogy/} & Profils, graphe de compétences, leçons, exercices, modèle de l'apprenant, recommandations. \\
\file{ui_api/}, \file{store/} & Protocole du panneau, proxy de synthèse vocale, page d'appairage ; persistance JSON/JSONL dans \file{~/.wysp/} avec masquage des secrets. \\
\bottomrule
\end{xltabular}

\subsubsection{La boucle de tour et les outils du LLM}

Un \textbf{tour} traite un message de l'utilisateur. Le cœur construit un contexte d'au plus 6\,000 jetons : prompt système générique, manifeste (types d'entités, catalogue d'actions), amorce du profil, contexte de l'adaptateur et observation du focus entourés de balises \code{<app\_data>} (contenu traité comme des données, jamais comme des instructions), état de l'apprenant, leçon en cours et conversation récente. Le LLM répond uniquement par des appels d'outils, au plus 8 tours d'outils par message.

\begin{xltabular}{\textwidth}{@{}L{3.0cm}Y L{2.3cm}@{}}
\toprule
\textbf{Outil} & \textbf{Effet} & \textbf{Capacité requise} \\
\midrule
\code{observe} & Résumé compact d'une observation & \code{observe} \\
\code{show\_entities} & Projection d'entités réelles sur le tableau (nœuds \pobs) & \code{observe} \\
\code{add\_concept}, \code{annotate} & Nœuds et annotations \pinf{} ou \pill, jamais \pobs & -- \\
\code{narrate} & Segments de narration avec repères (focus, surbrillance, étape) & -- \\
\code{highlight} & Mise en évidence dans l'application & \code{highlight} \\
\code{propose\_action} & Proposition à la porte de permission, mode \code{guide} ou \code{do} & \code{actions}, \code{execute} \\
\code{start\_exercise} & Démarrage d'un exercice du profil si ses liaisons se résolvent & \code{observe} \\
\code{lookup\_concept} & Entrée du glossaire du profil & -- \\
\code{ask\_clarification} & Question avec au plus 4 réponses proposées & -- \\
\bottomrule
\end{xltabular}

Les outils dont la capacité manque au manifeste ne sont pas proposés (propriété CORE-P10). Il n'existe aucun outil « exécuter » : c'est la garantie centrale du modèle de confiance.

\subsection{Modèle visuel et tableau}

Le tableau est un \textbf{modèle structuré}, pas une image : nœuds (avec \code{ref}, \code{kind}, provenance, confiance, propriétés à afficher), arêtes, étapes, annotations, focus et disposition. Le cœur publie des correctifs versionnés (\code{board.patch \{base\_version, version, ops\}}) ; le panneau les applique dans l'ordre et demande un instantané en cas de trou. La même fonction d'application existe en Python et en TypeScript, vérifiée par des fixtures communes (CORE-P3, PAN-P1).

\begin{itemize}[leftmargin=1.2em]
  \item \textbf{Projection.} Une entité devient un nœud d'identifiant \code{n:<ref>} ; les relations parent deviennent des arêtes \code{contains} ; les liens d'observation deviennent des arêtes typées. Le profil choisit les propriétés affichées par type (\code{board.props}). La projection est déterministe (CORE-P12), ce qui stabilise la disposition.
  \item \textbf{Validateur de provenance.} Tout nœud marqué \pobs{} dont la référence n'apparaît dans aucune observation de la session est rétrogradé en \pinf{} et journalisé (CORE-P1).
  \item \textbf{Dispositions.} Quatre suffisent pour les deux tuteurs et couvrent les exemples de la note de concept.
\end{itemize}

\begin{xltabular}{\textwidth}{@{}L{2.0cm}Y Y@{}}
\toprule
\textbf{Disposition} & \textbf{Blender} & \textbf{Web} \\
\midrule
\code{stack} & pile de modificateurs, ordre d'évaluation & ordre des filtres appliqués \\
\code{tree} & hiérarchie de la scène (collections, objets) & régions de la page, contrôles d'un formulaire \\
\code{graph} & objets, matériaux et références & dépendances entre filtres, KPI et tableau \\
\code{sequence} & étapes guidées d'une leçon & étapes de création d'un rapport \\
\bottomrule
\end{xltabular}

\begin{xltabular}{\textwidth}{@{}L{2.4cm}Y L{4.4cm}@{}}
\toprule
\textbf{Provenance} & \textbf{Style du nœud} & \textbf{Marqueur dans les sous-titres} \\
\midrule
\pobs & bordure pleine 2\,px, opacité complète, icône « œil » & aucun \\
\pinf & bordure pointillée, opacité 85\,\%, badge de confiance & « probablement » \\
\pill & fond hachuré, bordure grise, étiquette « exemple » & « par exemple » \\
\bottomrule
\end{xltabular}

\subsection{Voix et narration}

Nous retenons une chaîne en cascade plutôt qu'un modèle vocal de bout en bout : transcription dans le navigateur, raisonnement en texte dans le cœur, synthèse vocale par segments. Le texte reste ainsi la source de vérité : sous-titres, provenance et repères de synchronisation sont exacts, et la voix peut tomber en panne sans casser le tutorat.

\begin{figure}[h]
\centering
\begin{tikzpicture}[font=\sffamily\scriptsize, node distance=0.55cm,
  st/.style={draw=#1, fill=#1!8, rounded corners=2pt, align=center, minimum height=0.8cm, minimum width=1.9cm},
  arr/.style={-{Stealth[length=2mm]}, thick, color=surface}]
  \node[st=pb] (mic) {Micro\\\tiny appui pour parler};
  \node[st=pb, right=of mic] (stt) {Transcription\\\tiny Web Speech API};
  \node[st=pa, right=of stt] (core) {Cœur : tour\\\tiny LLM + outils};
  \node[st=pa, right=of core] (seg) {Segments\\\tiny $\leq$ 300 car. + repères};
  \node[st=pb, right=of seg] (tts) {Synthèse\\\tiny /api/tts ou navigateur};
  \node[st=pb, right=of tts] (spk) {Audio et\\sous-titres};
  \draw[arr] (mic) -- (stt); \draw[arr] (stt) -- (core); \draw[arr] (core) -- (seg); \draw[arr] (seg) -- (tts); \draw[arr] (tts) -- (spk);
  \draw[arr, color=danger, dashed] (spk.south) -- ++(0,-0.55) -| node[pos=0.25, below, text=danger] {interruption : parole détectée ou Échap $\rightarrow$ \code{user.interrupt}, arrêt en moins de 150\,ms} (mic.south);
\end{tikzpicture}
\caption{Chaîne vocale. En violet le panneau, en bleu le cœur.}
\end{figure}

\begin{itemize}[leftmargin=1.2em]
  \item \textbf{Appui pour parler} (barre d'espace ou bouton) par défaut : fiable en démonstration, pas de déclenchement intempestif.
  \item \textbf{Interruption} : pendant la narration, une détection d'énergie vocale (WebAudio, seuil calibré, annulation d'écho) ou la touche Échap arrête l'audio et envoie \code{user.interrupt}. Le cœur annule le tour et mémorise le dernier segment joué. « Continue » rejoue exactement les segments non entendus (CORE-P7, CORE-P8).
  \item \textbf{Clés} : la synthèse passe par \code{/api/tts} pour que la clé reste côté serveur ; la synthèse du navigateur sert de repli.
  \item \textbf{Latence visée} : premier son moins de 2,5\,s après la fin de la question, grâce au découpage en segments et à la prélecture du segment suivant.
\end{itemize}

\subsection{Pédagogie et vérification}

Le moteur pédagogique lit le \textbf{profil} de l'application : compétences avec prérequis, glossaire de concepts reliés à la documentation officielle, leçons en étapes (\code{explain}, \code{show}, \code{guide}, \code{do}, \code{check}) et exercices.

\begin{lstlisting}[language={}]
exercises:
  - id: bevel_segments
    prompt: {en: "Give the bevel 3 segments so the edge looks round."}
    skills: [bevel_modifier]
    bindings: {target: "$focus.object"}
    success: {all: [ ...predicates on the observation... ]}
    hints: [{en: "Look at the Modifier Properties tab (blue wrench)."},
            {en: "Find the Segments field in the Bevel panel."}]
    max_attempts: 3
\end{lstlisting}

\begin{itemize}[leftmargin=1.2em]
  \item \textbf{Exercices contextualisés} : un exercice n'est proposé que si ses liaisons se résolvent sur la scène ou la page réelle et si son critère n'est pas déjà satisfait (PED-P5, PED-P6).
  \item \textbf{Vérification} : à chaque changement d'origine utilisateur, le moteur évalue les prédicats sur une observation fraîche et enregistre une preuve (prédicat, valeur réelle, valeur attendue, identifiant d'observation, horodatage).
  \item \textbf{Trois niveaux de preuve} : \code{verified} (constaté dans l'application, poids 1,0), \code{answered} (réponse à un quiz, 0,6), \code{declared} (« c'est fait » non vérifiable, 0,3). Le panneau affiche toujours la différence. Les actions faites par Wysp ne comptent jamais (PED-P4).
  \item \textbf{Modèle de l'apprenant} : traçage bayésien des connaissances simplifié (glissement 0,1, chance 0,2, apprentissage 0,15), pondéré par le type de preuve et le nombre d'indices. Une compétence est acquise à 0,85. La profondeur des explications suit la maîtrise : bases sous 0,4, pratique guidée jusqu'à 0,85, défi au-delà.
\end{itemize}

\subsection{Confiance et sécurité}

\begin{figure}[h]
\centering
\begin{tikzpicture}[font=\sffamily\scriptsize,
  s/.style={draw=#1, fill=#1!8, rounded corners=2pt, align=center, minimum height=0.66cm, minimum width=1.75cm},
  arr/.style={-{Stealth[length=2mm]}, thick, color=surface},
  lbl/.style={font=\sffamily\tiny, text=muted, fill=white, inner sep=1pt}]
  \node[s=pa] (p) at (0,0) {proposée\\\tiny paramètres validés};
  \node[s=pa] (pol) at (2.4,0) {politique\\\tiny selon le risque};
  \node[s=okgreen] (auto) at (5.0,0.95) {approuvée\\\tiny auto ou autorisation};
  \node[s=pumpkin] (wait) at (5.0,-0.95) {en attente\\\tiny délai 60\,s};
  \node[s=muted] (deny) at (5.0,-2.35) {refusée\\\tiny ou expirée};
  \node[s=okgreen] (ex) at (7.7,0) {exécution\\\tiny adaptateur};
  \node[s=pa] (ver) at (10.1,0) {vérification\\\tiny post-conditions};
  \node[s=okgreen] (done) at (12.6,0.95) {réussie\\\tiny preuve enregistrée};
  \node[s=danger] (undo) at (12.6,-0.95) {échec\\\tiny annulation auto};
  \draw[arr] (p) -- (pol);
  \draw[arr] (pol) -- (auto);
  \draw[arr] (pol) -- (wait);
  \draw[arr] (auto) -- (ex);
  \draw[arr] (wait) -- node[lbl, pos=0.45, below right=-1pt] {clic ou « oui »} (ex);
  \draw[arr] (wait) -- node[lbl, right] {refus} (deny);
  \draw[arr] (ex) -- (ver);
  \draw[arr] (ver) -- (done);
  \draw[arr] (ver) -- (undo);
\end{tikzpicture}
\caption{Cycle de vie d'une proposition d'action.}
\end{figure}

\begin{itemize}[leftmargin=1.2em]
  \item \textbf{Provenance} : seule la projection d'observations crée des nœuds \pobs ; le validateur rattrape le reste. Au niveau d'intégration visuel, la confiance des entités est inférieure à 1 et affichée.
  \item \textbf{Porte de permission} : politique de la table des risques, autorisations limitées à une leçon pour les actions réversibles, clic obligatoire et aperçu pour les actions irréversibles, expiration à 60\,s. Une surcharge de politique par profil ne peut être que plus stricte (CORE-P4).
  \item \textbf{Vérifier ou annuler} : après exécution, nouvelle observation et évaluation des post-conditions ; en cas d'échec, annulation automatique et rapport avec preuve (CORE-P6). Chaque action Blender pousse aussi une étape d'annulation native, donc Ctrl+Z fonctionne.
  \item \textbf{Sécurité locale} : écoute sur \code{127.0.0.1} uniquement, contrôle de l'en-tête \code{Origin}, jeton d'appairage, aucune clé dans le navigateur, journaux masqués (CORE-P11).
  \item \textbf{Vie privée (web)} : extension active seulement sur les origines autorisées ; jamais de mot de passe, carte bancaire, code à usage unique, champ caché, cookie ou stockage dans les observations (WEB-P1, WEB-P6).
  \item \textbf{Injection de prompt} : le texte d'une page ou d'une scène est une donnée. Une page piégée peut au pire faire \emph{proposer} une action ; la porte et le catalogue fermé empêchent qu'elle s'exécute sans l'utilisateur.
  \item \textbf{Audit} : propositions, décisions (qui, par quel canal), exécutions, vérifications et annulations dans \file{~/.wysp/audit/*.jsonl} (CORE-P5).
\end{itemize}
